\chapter{Why good documents still fail}
\label{ch:problem}

The author sees the repository, the model's history, and the reason each
qualification was added. The reader sees a rendered document for the first
time. In the failed funding proposal that difference was enough to turn a
correct draft into a confusing one. Humanvoice begins with that change in
vantage point.

\section{The first five minutes}

An expert reader does not begin by grading sentence style. The reader tries to
build a small working model of the document:

\begin{enumerate}
\item What decision or problem is this document about?
\item What is the proposed object, mechanism, or intervention?
\item Why is the obvious alternative insufficient?
\item Which evidence would change my view?
\item What action is the author asking me to take?
\end{enumerate}

If those answers are not available, later polish has diminishing value. The
reader must spend scarce attention discovering the document's structure before
deciding whether the underlying work deserves attention. In the internal case,
the first pages named phases, files, and review states before naming the action
that the funder was meant to approve. Each local sentence was defensible;
the sequence was not.

This is why ordinary readability formulas are insufficient. A short sentence
can conceal an undefined mechanism, and a long sentence can be efficient when
its actors and quantities are familiar. Readability is a relation between a
document, a task, and a reader. It cannot be inferred from a single number.

Figure~\ref{fig:reader-burden-cycle} makes the economic problem visible. The
cost arises at the missing relation, where the reader must supply work the
document should have done.

\hvFigureReaderBurden

\section{How a correct draft becomes difficult to read}

The internal cases suggest a sequence of failures. The sequence is useful not
as a universal taxonomy but as a way to decide where a product intervention can
help.

\begin{figure}[H]
\centering
\resizebox{\textwidth}{!}{%
\begin{tikzpicture}[node distance=4mm and 5mm,font=\small]
  \node[hvstage,text width=32mm] (source) {Correct technical\\source};
  \node[hvwarn,text width=32mm,right=of source] (register) {Private register\\leaks into prose};
  \node[hvwarn,text width=32mm,right=of register] (rhythm) {Repeated shapes\\mask the point};
  \node[hvwarn,text width=32mm,right=of rhythm] (defense) {Qualifications\\replace reasons};
  \node[hvwarn,text width=32mm,right=of defense] (substance) {Mechanism or\\evidence stays implicit};
  \node[hvgood,text width=34mm,below=13mm of substance] (reader) {Reader cannot recover\\purpose or action};
  \draw[hvflow] (source)--(register)--(rhythm)--(defense)--(substance);
  \draw[hvflow] (substance)--(reader);
  \draw[hvflow] (register.south) to[bend right=18] (reader.west);
\end{tikzpicture}
}%
\caption{A recurring path from local correctness to reader failure. The arrows
are diagnostic opportunities, not a claim that every document follows all five
steps.}
\label{fig:failure-layers}
\end{figure}

\paragraph{Private register.} A project identifier, phase label, file path, or
agent status can be meaningful to the authors and opaque to a new reader. The
problem is not that the term is technical. The problem is that its referent and
importance are unavailable in the document.

\paragraph{Mechanical regularity.} A document assembled through repeated
templates may open every section in the same way, use the same transition, or
close every discussion with the same kind of qualification. These patterns are
observable, but their presence is not itself a verdict. They matter when they
make the reader predict the form instead of following the argument.

\paragraph{Defensive prose.} A sentence can list what the project does not
claim, what a tool cannot guarantee, or what a reviewer must not infer without
giving the reader the positive mechanism. A boundary is useful when it prevents
a plausible mistake. It is a burden when it substitutes for an explanation.

\paragraph{Missing substance.} A document may name a result without showing
the mechanism connecting evidence to conclusion, or name a feature without
showing the user decision it changes. No surface diagnostic can establish that
connection by itself.

\paragraph{Reader failure.} The final test is whether the intended reader can
reconstruct the document's purpose and make the requested judgment. This is not
the same as agreement. A reader can understand a proposal and reject it for
good reasons. Humanvoice must preserve that distinction.

\section{Why the incumbent workflow stops too early}

Today a linter reports a phrase, a model proposes a rewrite, version control
records a file change, and a human reader makes the final judgment. The links
between those activities are usually informal: a reviewer pastes a comment into
a ticket, an author edits a paragraph, and a later reader is asked to trust
that nothing important changed. That is the structural blind spot.

Three costs follow.

First, diagnosis is not located. A comment such as ``the argument is hard to
follow'' may be accurate but gives the author no shared object on which to act.
Second, revision is not protected. A fluent rewrite can alter a sign, a
denominator, a citation, or the scope of a conclusion without an obvious visual
signal. Third, acceptance is not recorded as an outcome. The team remembers
that a senior person ``looked at it,'' but cannot compare reader time or the
number of rounds across documents.

Humanvoice earns its place only if it reduces those three costs together. A new
dashboard that adds another queue of findings makes the problem worse. A
protected diff that no author can understand makes safety too expensive. A
reader questionnaire detached from the actual decision produces a number
without a use.

\section{The people and the decision}

The product has four human roles, even when one person temporarily occupies
more than one role.

\begin{description}
\item[Author or editor.] Owns the claims, evidence, wording, and decision to
  accept or reject a suggested repair.
\item[Technical owner.] Supplies the document's format, protected-object rules,
  privacy boundary, and integration point with the repository or editor.
\item[Reader.] Represents the person whose judgment the document is meant to
  support. The reader may be a funder, executive, reviewer, or collaborator.
\item[Sponsor.] Decides whether the saved attention and reduced risk justify
  further investment. The sponsor may be the technical owner, but need not be.
\end{description}

The author needs actionable findings and safe comparison. The reader needs a
document that does not require privileged project context. The sponsor needs a
credible comparison with the incumbent bundle. These are compatible needs, but
they are not the same metric.

\section{A useful product boundary}

The product should be judged at the boundary where it changes a real review
decision. It is not a general-purpose writing environment, a style-imposition
service, or an authorship detector. It is a review path beside the author's
existing editor and repository.

\begin{takeawaybox}{The boundary that matters}
Humanvoice owns the evidence trail between a source snapshot and a reader
decision. The author owns the prose. The reader owns acceptance. The sponsor
owns the decision to continue.
\end{takeawaybox}

This boundary also explains why a local-first design is more than a privacy
preference. Technical proposals can contain unpublished results, client data,
or decisions that should not be sent to an external model by default. A system
that cannot show where a finding came from cannot offer a useful privacy or
meaning guarantee either. Source locations, protected spans, and versioned
review decisions are the minimum common language between the four roles. The
next chapter turns from the boundary to the investment question: what burden
would a bounded build have to remove, and what evidence would make the result
worth funding?
